\begin{table}[H]
  \centering
  \caption{Matrices de transición de pobreza, monetaria por ingreso vs.
  IPM (\% de fila, panel emparejado 1 a 1). Estado inicial en filas,
  estado final en columnas.}
  \label{tab:matriz_transicion}
  \footnotesize
  \begin{tabular}{ll cc cc}
    \toprule
    & & \multicolumn{2}{c}{\textbf{Monetaria}} & \multicolumn{2}{c}{\textbf{IPM}} \\
    \textbf{Periodo} & \textbf{Estado inicial} & \textbf{No pobre} & \textbf{Pobre} & \textbf{No pobre} & \textbf{Pobre} \\
    \midrule
    \multirow{2}{*}{2010 $\rightarrow$ 2013} & No pobre & 76.6\% & 23.4\% & 87.9\% & 12.1\% \\
     & Pobre & 29.7\% & 70.3\% & 54.5\% & 45.5\% \\
    \addlinespace
    \multirow{2}{*}{2013 $\rightarrow$ 2016} & No pobre & 80.1\% & 19.9\% & 90.3\% & 9.7\% \\
     & Pobre & 38.2\% & 61.8\% & 52.3\% & 47.7\% \\
    \bottomrule
  \end{tabular}
  \begin{minipage}{0.95\textwidth}
    \vspace{4pt}
    \footnotesize \textit{Nota:} $n=8{,}218$ (2010$\to$2013) y $n=6{,}911$
    (2013$\to$2016) hogares con emparejamiento 1 a 1 -- idéntico universo para
    monetaria e IPM (mismo panel, solo cambia la clasificación de pobreza);
    excluye 511 (2010--2013) y 1{,}190 (2013--2016) casos de hogares que se
    dividieron entre olas. Fuente: cálculos propios con base en ELCA 2010,
    2013 y 2016.
  \end{minipage}
\end{table}